% ============================================================================
%  CAPÍTULO I: INTRODUCCIÓN, CLIENTE Y ALCANCE
% ============================================================================

\section{Resumen ejecutivo}

La Cooperativa de Ahorro y Crédito Horizonte PyME S.R.L. gestiona sus créditos
y ahorros sobre una base de datos MySQL 5.7, motor que perdió el soporte
oficial en octubre de 2023. Para resolver ese riesgo se migró la base a
PostgreSQL 16, en un entorno desplegado íntegramente sobre contenedores Docker
y orquestado con Docker Compose. La migración incluyó el diagnóstico del
origen, el rediseño del esquema con tipos nativos, un proceso de carga
reproducible, el endurecimiento (\emph{hardening}) del servidor destino, roles
de mínimo privilegio, enmascaramiento de datos personales, respaldo cifrado y
pruebas automatizadas de cada resultado.

La migración se ejecutó sobre un conjunto de datos realista y determinista:
8 tablas, 2 vistas, 19 entradas de índice, 55 socios, 66 cuentas de ahorro,
50 solicitudes de crédito, 32 préstamos, 516 cuotas de amortización y 740
transacciones, para un total de 1\,469 filas migradas. El origen ocupa
0.28 MB y presenta incompatibilidades deliberadas con PostgreSQL: 7 columnas
\texttt{ENUM}, 10 columnas \texttt{DATETIME}, 4 columnas \texttt{TINYINT(1)},
claves \texttt{INT UNSIGNED AUTO\_INCREMENT}, collation \texttt{utf8\_general\_ci}
y ausencia total de claves foráneas. Además, el usuario de aplicación
\texttt{app\_legacy\_user} conserva un \texttt{GRANT ALL} sobre toda la base,
lo que constituye un hallazgo de auditoría de primer orden.

El destino se diseñó con \texttt{BOOLEAN}, \texttt{TIMESTAMPTZ},
identidades \texttt{BIGINT GENERATED ALWAYS AS IDENTITY}, restricciones
\texttt{CHECK}, 9 claves foráneas y 10 índices orientados a los reportes de
cartera y saldos. El endurecimiento incluye TLS/SSL obligatorio con
\texttt{hostssl}, autenticación \texttt{scram-sha-256}, revocación de
privilegios de \texttt{PUBLIC}, tres roles granulares (\texttt{admin\_dba},
\texttt{app\_core\_rw} y \texttt{auditor\_ro}), una vista enmascarada
\texttt{vw\_socios\_audit} y registro de conexiones, desconexiones y DDL.
El respaldo del destino se genera con \texttt{pg\_dump -Fc} y se cifra con
AES-256-CBC y PBKDF2; la restauración se probó y devolvió los 55 socios
esperados.

La verificación final arrojó 75 pruebas automatizadas en estado \texttt{PASS},
ninguna en \texttt{FAIL}, cero huérfanos en las 9 relaciones, conteos idénticos
entre origen y destino en las 8 tablas, rechazo efectivo de conexiones sin SSL
y acceso remoto directo de \texttt{postgres} bloqueado. Las consultas críticas
se ejecutan en el servidor PostgreSQL en 0.236 ms y 0.317 ms, muy por debajo de
cualquier umbral operativo para el volumen analizado. El proyecto se entrega
con scripts, evidencias y anexos.

\section{Descripción del cliente}

La Cooperativa de Ahorro y Crédito Horizonte PyME S.R.L. es una institución
financiera de alcance regional orientada a microcréditos para emprendedores,
comercios locales y pequeñas y medianas empresas. Su operación se organiza en
cinco sucursales (Central, Norte, Sur, Este y Oeste) y gira en torno a tres
procesos de negocio: la captación de ahorro, la colocación de crédito y el
seguimiento de la cartera. Los productos de crédito vigentes son personal,
empresarial, educativo, médico y agrícola, cada uno con montos máximos, tasas
y plazos diferenciados.

El sistema de información administra el ciclo de vida completo del socio:
registro y actualización de datos personales, apertura de cuentas de ahorro,
presentación y evaluación de solicitudes de crédito, desembolso, generación
de la tabla de amortización, registro de pagos y control de morosidad. Los
reportes principales son la cartera en mora (\texttt{vw\_cartera\_mora}) y el
resumen de saldos por socio (\texttt{vw\_resumen\_saldos\_socio}). La
información tratada es financiera y personal: nombres, documentos de identidad,
teléfonos, correos, direcciones, números de cuenta, saldos, montos de crédito
y historial de pagos.

Los interesados (\emph{stakeholders}) del proyecto son la Gerencia General,
que necesita continuidad operativa y información oportuna para la toma de
decisiones; la Unidad de Riesgos, que depende de la integridad de la cartera;
la Unidad de Auditoría, que requiere trazabilidad y acceso controlado; el
Área de Tecnología de la Información, responsable de la operación; y los
socios, cuyos datos deben permanecer confidenciales y disponibles. La
migración debe preservar la confianza de todos ellos: un dato perdido, una
cuenta mal conciliada o una filtración de identidad tendrían consecuencias
reputacionales, regulatorias y económicas.

\subsection{Requisitos del cliente}

A partir de las necesidades expresadas por los interesados se derivaron los
requisitos que el sistema destino debía satisfacer. La
Tabla~\ref{tab:requisitos-funcionales} enumera los requisitos funcionales.

\begin{table}[H]
  \centering
  \caption{Requisitos funcionales del cliente}
  \label{tab:requisitos-funcionales}
  \contenidoConFuente{%
    \begin{tablaAPA}
      \begin{tabular}{@{}C{0.9cm}L{4.4cm}L{10.2cm}@{}}
        \toprule
        \textbf{N.} & \textbf{Requisito} & \textbf{Descripción} \\
        \midrule
        1 & Migración íntegra & Traslado de las 8 tablas y las 2 vistas a PostgreSQL 16 preservando identificadores, integridad referencial y zona horaria. \\
        2 & Endurecimiento del destino & TLS obligatorio, autenticación \texttt{scram-sha-256}, roles de mínimo privilegio y registro de conexiones y cambios de estructura. \\
        3 & Enmascaramiento & Vista de auditoría con CI/NIT, teléfono, correo y cuenta enmascarados, sin exponer las tablas base. \\
        4 & Respaldo cifrado y restauración & Generación de un respaldo cifrado y prueba de su restauración completa. \\
        5 & Automatización y evidencias & Auditoría, migración y verificación con scripts reproducibles y evidencias reales. \\
        \bottomrule
      \end{tabular}
    \end{tablaAPA}%
  }{Elaboración propia.}
\end{table}

La Tabla~\ref{tab:requisitos-no-funcionales} complementa la lista anterior con
los requisitos no funcionales.

\begin{table}[H]
  \centering
  \caption{Requisitos no funcionales del cliente}
  \label{tab:requisitos-no-funcionales}
  \contenidoConFuente{%
    \begin{tablaAPA}
      \begin{tabular}{@{}C{0.9cm}L{4.4cm}L{10.2cm}@{}}
        \toprule
        \textbf{N.} & \textbf{Requisito} & \textbf{Descripción} \\
        \midrule
        1 & Infraestructura & Ejecución de todo el entorno con Docker Compose, sin instalar motores en el sistema anfitrión. \\
        2 & Sistema operativo y red & Linux en el laboratorio y puertos publicados solo en \texttt{127.0.0.1}. \\
        3 & Zona horaria & \texttt{America/La\_Paz} para fechas y reportes. \\
        4 & Entregables & Informe con la plantilla UPDS y demostración en vivo de 15--20 minutos, sin modificar \texttt{plantilla-latex/}. \\
        5 & Datos & Datos sintéticos, sin información personal real. \\
        \bottomrule
      \end{tabular}
    \end{tablaAPA}%
  }{Elaboración propia.}
\end{table}

Cada requisito se verifica en el Capítulo~IV y se resume en la lista de
cotejo del Anexo~C. La combinación de TLS, roles de mínimo privilegio,
enmascaramiento y respaldo cifrado exigida por el cliente se corresponde con
los controles recomendados por CIS, OWASP y NIST
\citep{cis2024postgres,owasp2023asvs,nist80053}.

\section{Problema de contexto}

MySQL 5.7 alcanzó el fin de su vida útil el 31 de octubre de 2023
\citep{oracle2023mysql57}. A partir de esa fecha Oracle dejó de publicar
parches de seguridad y correcciones para la rama 5.7, de modo que cualquier
vulnerabilidad descubierta queda sin remedio oficial. Para una entidad
financiera, operar sobre un motor sin soporte implica aceptar un riesgo
tecnológico y regulatorio que crece con el tiempo. El problema no es
únicamente la versión: el esquema heredado arrastra decisiones de diseño que
hoy resultan costosas.

En primer lugar, el usuario de aplicación \texttt{app\_legacy\_user} tiene
\texttt{GRANT ALL PRIVILEGES} sobre toda la base. Si la aplicación sufre una
inyección SQL o un compromiso de credenciales, el atacante hereda permisos de
lectura, escritura, eliminación y modificación de estructura. No existe
separación entre la cuenta que opera el sistema y la cuenta que administra la
base, ni entre la operación y la auditoría. En segundo lugar, las conexiones
no exigen cifrado: el tráfico entre la aplicación y el motor puede viajar en
texto claro dentro de la red, lo que expone credenciales y datos financieros.
En tercer lugar, el esquema no declara claves foráneas, de modo que la
integridad referencial depende por completo del código de la aplicación; una
carga manual, un \emph{script} defectuoso o un borrado accidental pueden dejar
cuentas sin socio, cuotas sin préstamo o transacciones sin cuenta.

A lo anterior se suman las incompatibilidades de tipos. \texttt{TINYINT(1)}
se usa como booleano, pero MySQL acepta valores distintos de 0 y 1;
\texttt{ENUM} restringe los estados en el motor, pero no se traduce de forma
directa a PostgreSQL; \texttt{DATETIME} no almacena zona horaria, lo que
introduce ambigüedad en operaciones distribuidas; \texttt{INT UNSIGNED}
duplica el rango positivo pero no existe como tipo nativo en PostgreSQL; y la
collation \texttt{utf8\_general\_ci} corresponde a \texttt{utf8mb3}, una
codificación de tres bytes que no cubre todos los caracteres Unicode. La
Tabla~\ref{tab:riesgos-origen} resume los riesgos identificados y su efecto
potencial.

\begin{table}[H]
  \centering
  \caption{Riesgos detectados en el sistema de origen}
  \label{tab:riesgos-origen}
  \contenidoConFuente{%
    \begin{tablaAPA}
      \begin{tabular}{@{}L{3.1cm}L{5.3cm}L{6.2cm}@{}}
        \toprule
        \textbf{Riesgo} & \textbf{Evidencia} & \textbf{Impacto potencial} \\
        \midrule
        Motor sin soporte & MySQL 5.7.44, EOL 31/10/2023 & Vulnerabilidades sin parche; incumplimiento de políticas de seguridad. \\
        Privilegios excesivos & \texttt{GRANT ALL} a \texttt{app\_legacy\_user} & Compromiso total de la base ante una inyección o robo de credenciales. \\
        Tráfico sin cifrado & Sin TLS obligatorio en el origen & Exposición de credenciales, saldos e identidades en la red. \\
        Ausencia de claves foráneas & 0 \texttt{FOREIGN KEY} en el esquema & Huérfanos, inconsistencias y reportes financieros erróneos. \\
        Tipos ambiguos & 7 \texttt{ENUM}, 10 \texttt{DATETIME}, 4 \texttt{TINYINT(1)} & Conversiones silenciosas, pérdida de zona horaria y validaciones débiles. \\
        Collation limitada & \texttt{utf8\_general\_ci} (\texttt{utf8mb3}) & Caracteres no representables y comparaciones inconsistentes. \\
        Auditoría insuficiente & Sin registro de conexiones ni DDL & Imposibilidad de reconstruir accesos y cambios. \\
        Respaldo no cifrado & \texttt{mysqldump} en texto claro & Exposición de datos si el archivo de respaldo se filtra. \\
        \bottomrule
      \end{tabular}
    \end{tablaAPA}%
  }{Elaboración propia con datos de \texttt{evidence/01\_inventario\_origen.txt}.}
\end{table}

La cooperativa necesita modernizar su plataforma de datos sin detener la
operación, sin perder información y sin repetir los defectos de seguridad del
sistema heredado. La migración debe ser verificable, reversible y reproducible
en cualquier equipo del grupo de trabajo.

\section{Solución propuesta}

La solución se apoya en cinco decisiones técnicas. La migración a
PostgreSQL 16 aporta un motor con soporte activo, cumplimiento ACID, tipos
nativos para booleanos, marcas de tiempo con zona horaria, identidades,
restricciones \texttt{CHECK}, claves foráneas e índices avanzados
\citep{postgresql2024}. El entorno completo se ejecuta con Docker Compose: un
servicio MySQL 5.7 para el origen, un servicio PostgreSQL 16 para el destino y
un contenedor migrador con Python, \texttt{pymysql} y \texttt{psycopg2}. Así se
evita instalar motores en el sistema anfitrión y se garantiza que cualquier
integrante levante el mismo escenario con un solo comando
\citep{docker2024,compose2024}.

Las transformaciones se definieron de forma explícita, regla por regla, en
lugar de depender de una conversión automática genérica. El volumen es pequeño
(0.28 MB), pero las reglas son específicas: convertir \texttt{TINYINT(1)} a
\texttt{BOOLEAN}, \texttt{ENUM} a \texttt{VARCHAR} con \texttt{CHECK},
\texttt{DATETIME} a \texttt{TIMESTAMPTZ} en la zona \texttt{America/La\_Paz},
preservar los identificadores originales con \texttt{OVERRIDING SYSTEM VALUE}
y reajustar las secuencias de identidad. Un ETL explícito permite auditar cada
regla y repetir la carga de forma determinista.

El proceso de migración se construyó sobre \textbf{pgloader}
con schema de staging (\texttt{stg}), validación SQL previa (\texttt{precheck}),
carga en una sola transacción con rollback garantizado (\texttt{postcheck}),
reajuste de secuencias y \texttt{ANALYZE}. El ETL artesanal en Python
(\texttt{src/migration/migrate.py}) fue archivado como legado y ya no se
ejecuta.

La seguridad se aplicó por capas. El perímetro exige TLS;
la autenticación usa \texttt{scram-sha-256}; la autorización se basa en roles
con privilegios mínimos; los datos sensibles se exponen a auditoría solo a
través de una vista enmascarada; las conexiones y el DDL se registran; y los
respaldos se cifran con AES-256-CBC. La Tabla~\ref{tab:solucion-componentes}
relaciona cada componente con su propósito y su evidencia.

\begin{table}[H]
  \centering
  \caption{Componentes de la solución propuesta}
  \label{tab:solucion-componentes}
  \contenidoConFuente{%
    \begin{tablaAPA}
      \begin{tabular}{@{}L{3.5cm}L{6.0cm}L{5.0cm}@{}}
        \toprule
        \textbf{Componente} & \textbf{Propósito} & \textbf{Evidencia} \\
        \midrule
        \texttt{docker-compose.yml} & Orquestar origen, destino, migrador y documentación & Servicios \emph{healthy}; puertos en \texttt{127.0.0.1}. \\
        \texttt{src/origin/} & Esquema y datos legacy reproducibles & 8 tablas, 2 vistas, 1\,469 filas. \\
        \texttt{src/target/} & Esquema destino, roles y configuración endurecida & 9 claves foráneas, 3 roles, SSL y \texttt{pg\_hba}. \\
        \texttt{src/migration/} & pgloader + staging, precheck, postcheck y rollback transaccional & 1\,469 filas, 0 rechazos, 0 errores. \\
        \texttt{scripts/} & Auditoría, migración, respaldo y pruebas & 3 scripts y 3 evidencias reales. \\
        \texttt{docs/} & Informe, anexos y guion de demostración & PDF compilable y guion de 15--20 min. \\
        \bottomrule
      \end{tabular}
    \end{tablaAPA}%
  }{Elaboración propia.}
\end{table}

\section{Alineación con los ODS 8 y 9}

El proyecto se alinea con el Objetivo de Desarrollo Sostenible 8 (Trabajo
decente y crecimiento económico) porque mejora la capacidad de la cooperativa
para otorgar microcréditos de forma oportuna. La meta 8.3 promueve la
formalización y el crecimiento de las micro, pequeñas y medianas empresas; la
meta 8.10 busca fortalecer la capacidad de las instituciones financieras para
ampliar el acceso a servicios bancarios y financieros \citep{un2023ods8}. Al
reducir el tiempo de aprobación y desembolso, y al disminuir el riesgo de
errores en la cartera, la migración contribuye directamente a que más
emprendedores accedan a capital de trabajo.

El ODS 9 (Industria, innovación e infraestructura) se relaciona con la
modernización de una infraestructura crítica de datos. La
meta 9.1 promueve infraestructuras fiables, sostenibles y resilientes; la meta
9.5 fortalece la investigación y la capacidad tecnológica; y la meta 9.c
busca ampliar el acceso a las tecnologías de la información y las
comunicaciones \citep{un2023ods9}. La contenerización, el cifrado, la
automatización de pruebas y la documentación reproducible son aportes
concretos a una infraestructura más resiliente y gobernable.

\begin{table}[H]
  \centering
  \caption{Contribución del proyecto a los ODS 8 y 9}
  \label{tab:ods}
  \contenidoConFuente{%
    \begin{tablaAPA}
      \begin{tabular}{@{}L{1.8cm}L{2.2cm}L{5.6cm}L{4.9cm}@{}}
        \toprule
        \textbf{ODS} & \textbf{Meta} & \textbf{Aporte del proyecto} & \textbf{Indicador de logro} \\
        \midrule
        ODS 8 & 8.3 & Agiliza la evaluación y el desembolso de microcréditos a pymes y emprendedores. & Tiempo de consulta de cartera y saldos en el orden de milisegundos. \\
        ODS 8 & 8.10 & Amplía el acceso a servicios financieros con información íntegra y disponible. & 55 socios y 66 cuentas migrados sin pérdida. \\
        ODS 9 & 9.1 & Sustituye una infraestructura obsoleta por una plataforma resiliente y reproducible. & Entorno Docker Compose levantado desde cero. \\
        ODS 9 & 9.5 & Incorpora automatización, pruebas y documentación técnica reutilizable. & 75 pruebas PASS y 3 evidencias reales. \\
        ODS 9 & 9.c & Facilita el acceso seguro a TIC mediante TLS, SCRAM y roles mínimos. & Conexiones sin SSL rechazadas; 7 roles auditados. \\
        \bottomrule
      \end{tabular}
    \end{tablaAPA}%
  }{Elaboración propia con base en \citet{un2023ods8} y \citet{un2023ods9}.}
\end{table}


